\documentclass[../main.tex]{subfiles}
\graphicspath{{\subfix{../images/}}}

\begin{document}
% \newcommand{\casedaysabs}{13,207,978}
\newcommand{\casedaysabssevcvd}{12,549}
\newcommand{\casedaysabssevresp}{18,159}
\newcommand{\casedayshyb}{13,207,978}
\newcommand{\cazipcodes}{1,721}
\newcommand{\controldaysabs}{44,819,948}
\newcommand{\controldayshyb}{44,819,948}
\newcommand{\copdabscihigh}{1.4}
\newcommand{\copdabscilow}{1.1}
\newcommand{\copdabsor}{1.2}
\newcommand{\copdhybcihigh}{1.5}
\newcommand{\copdhybcilow}{1.2}
\newcommand{\copdhybor}{1.3}
\newcommand{\cvdcasedayssev}{12,549}
\newcommand{\highestencountersexplevel}{Severe}
\newcommand{\meanwfpm}{0.2}
\newcommand{\medianduration}{85.8}
\newcommand{\mostcommonsevabs}{Severe}
\newcommand{\mostcommonsevhyb}{Mild}
\newcommand{\npspsevents}{64}
\newcommand{\psychcasedayssev}{5,099}
\newcommand{\respabscihigh}{1.1}
\newcommand{\respabscilow}{1.0}
\newcommand{\respabsor}{1.0}
\newcommand{\respcasedayssev}{18,159}
\newcommand{\resphybcihigh}{1.1}
\newcommand{\resphybcilow}{1.0}
\newcommand{\resphybor}{1.1}
\newcommand{\respintabscihigh}{1.2}
\newcommand{\respintabscilow}{1.0}
\newcommand{\respintabsor}{1.1}
\newcommand{\respinthybcihigh}{1.2}
\newcommand{\respinthybcilow}{1.0}
\newcommand{\respinthybor}{1.1}
\newcommand{\secondmostcommonsevabs}{Mild}
\newcommand{\secondmostcommonsevhyb}{Moderate}
\newcommand{\zipcodedayswf}{1,229,438}
\newcommand{\zipcodeswithevents}{876}
\newcommand{\zipdays}{11,527}
\newcommand{\zipdaysdualexp}{7,656}
\newcommand{\zipevents}{13.2}


\textbf{Online Data Supplement.}
\\
\\
\textbf{Authors:}
\\
Lauren Blair Wilner
Caitlin G. Jones-Ngo
Bethany Kwoka
Vivian Do
Chen Chen
Joel D. Kaufman
Tarik Benmahrnia
Joan A. Casey

\pagebreak
\section{Supplemental methods.}
\subsection{Circuit-level data description.}
The PSPS dataset contains circuit-level information for each PSPS event across California's three major utilities (Pacific Gas and Electric Company [PG\&E], San Diego Gas \& Electric [SDG\&E], and Southern California Edison [SCE]). A circuit is a network of electrical lines that delivers electricity from a power source to a defined service area of customers. A PSPS event may impact multiple circuits, so each circuit-event had a start time, end time, duration, and number of residential and commercial customers impacted. The circuits were provided as geospatial polygons with an 18-meter buffer around the circuit polyline that was approximated to include all customers potentially served by that circuit. The PSPS dataset contained the number of customers affected by each outage along the circuit. The data contained information on both commercial and residential customers. 

We overlaid these circuits with high-resolution gridded estimates of population from the 2020 census in California to estimate (1) the number of customers theoretically served by the circuit, and (2) the number of customers affected and the number of people served at the zip code level. In the absence of more detailed information on customer distribution, we assumed an even distribution of customers across the circuit and performed two spatial overlays of the circuits with the gridded population data. The first used the number of customers impacted on a circuit and proportionally split them between zip codes. The second used the population pixels that intersected with each circuit to determine the population served by the circuit. By summing each of these metrics by zip code, we determined the number of customers impacted and the population served on each circuit for each event. Notably, the denominator remained constant from event to event, while the numerator changed.

\subsection{Absolute (primary) and relative (secondary) PSPS exposure metrics.}
We derived and included two daily-zip code exposure metrics meant to capture different aspects of PSPS events. The first metric was the total number of customers impacted in each zip code. Because of highly variable population counts within zip codes, we also constructed a secondary relative metric meant to measure the severity of the PSPS event. The secondary metric was computed by multiplying the total number of customers without power by the proportion of customers affected in each zip code. This proportion was calculated using the numerator and denominator described in the \textit{Circuit-level data description}. Exposure for both metrics was categorized as no exposure or mild, moderate, or severe exposure based on rounded tertiles of the non-zero exposure values of each metric for each event-zip code. For our primary metric, the magnitude of same-day PSPS events were classified as mild when $\leq$ 61 customers were without power, moderate when the number of customers without power was $>$ 61 and $\leq$ 593, and severe when $>$ 593 customers were without power. For our secondary metric, magnitude for same-day events was again classified as mild (when our relative metric was $\leq$ 1), moderate (1 $<$ relative metric $\geq$ 145), or severe (when our relative metric was $>$ 145). For our primary metric, the magnitude of \text{lag}$_{0\text{-}3}$ PSPS events were classified as mild when $\leq$ 45 customers were without power, moderate when the number of customers without power was $>$ 45 and $\leq$ 458, and severe when $>$ 458 customers were without power. For our secondary metric, magnitude for \text{lag}$_{0\text{-}3}$ events was again classified as mild (when our relative metric was $\leq$ 1), moderate (1 $<$ relative metric $\geq$ 139), or severe (when our relative metric was $>$ 139). These tertiles were calculated at the event-zip code level to ensure equal weighting of longer and shorter events and retain the basis of the severity evaluation being on number of people rather than duration.  

Given the variability in zip code size, these absolute and relative metrics enabled us to measure exposure to PSPS outages, both in terms of the actual number of customers affected (absolute) and in terms of the portion of customers affected (relative). We leveraged the absolute metric for the primary analysis, and used our relative metric for a secondary.

The PSPS dataset included customers rather than individual people impacted. Customers impacted included both commercial and residential customers, meaning businesses and residences (e.g., individual apartments, single-family residences). This means that our primary exposure metric reflected the number of customers affected, not the number of people affected. For our secondary relative metric, the numerator included customers impacted, while the denominator relied on gridded population data representing individual people. This meant we consistently underestimated the proportion of people without power, but we do not expect this to be differential with respect to outcomes. 

\subsection{Model equation.}

\textbf{Regression Equation}

\begin{equation}
\log\left(\frac{p_{ij}}{1-p_{ij}}\right) = \alpha_i + \beta_1 \cdot PSPS_{ij} + \beta_2 \cdot PM_{2.5,ij} + \beta_3 \cdot (PSPS_{ij} \times PM_{2.5,ij}) + f(Temp_{max,ij})
\end{equation}

where: $p_{ij} = P(Y_{ij} = 1 | X_{ij}, C_{ij})$ for a case $(j = 1)$ or a control $(j = 0)$ in stratum $ij$. $i$ represents each stratum and $j$ represents a case or control in that stratum.

\vspace{1em}

\textbf{Parameter Definitions}

\begin{itemize}
\item $\alpha_i$: stratum-specific intercept that accounts for the baseline risk of the health outcome in each stratum, or matched set in this case. 
\item $\beta_1$: log odds ratio for emergency department visit or hospital admission associated with exposure to a PSPS event (binary indicator), assuming no wildfire PM$_{2.5}$ exposure and adjusted for temperature
\item $\beta_2$: log odds ratio for emergency department visit or hospital admission per \SI[per-mode=symbol]{10}{\micro\gram\per\cubic\metre} higher 5-day average wildfire PM$_{2.5}$ concentration, assuming no PSPS event exposure and adjusted for temperature
\item $\beta_3$: log odds ratio for the multiplicative interaction between PSPS event exposure and wildfire \PM (per \SI[per-mode=symbol]{10}{\micro\gram\per\cubic\metre}). This is the additional effect of \SI[per-mode=symbol]{10}{\micro\gram\per\cubic\metre} higher wildfire \PM in the presence of a PSPS event. 
\item $f(\mathrm{Temp_{max}})$: natural cubic spline with 3 degrees of freedom for daily maximum temperature.
\end{itemize}

\vspace{1em}

The joint effect of a PSPS event and \SI[per-mode=symbol]{10}{\micro\gram\per\cubic\metre} higher \PM exposure was calculated as: $\beta_1 + \beta_2 + \beta_3$. 

\subsection{Sensitivity analysis using outage durations.}
While we omitted duration from our primary analysis, longer PSPS events are likely to have greater health impacts. To assess this, we conducted a sensitivity analysis where we used number of hours without power in the prior week, irrespective of the number of affected households, as our primary PSPS exposure.

To calculate daily PSPS outage duration by zip code, we expanded each PSPS event to hourly intervals based on its start and end times. We rounded the start time down and the end time up to the nearest hour. We then de-duplicated overlapping hours across events within the same zip code to avoid double-counting hours when multiple PSPS events affected the same zip code during the same hour. Finally, we aggregated the unique outage hours to the daily level in order to generate the number of hours out per zip code-day (maximum 24 hours per day due to no double-counting hours). We then summed the number of hours without power in the week prior to index or referent day and used this as the PSPS exposure in our models.


\pagebreak
\section{Supplemental tables.}
\begin{table}[h]
\centering
\begin{tabular}{|p{0.18\textwidth}|p{0.39\textwidth}|p{0.39\textwidth}|}
\hline
\textbf{Health endpoint} & \textbf{\textit{ICD-9 Codes}} & \textbf{\textit{ICD-10 Codes}} \\
\hline
Respiratory & 460-466, 471, 472, 477, 478, 480-487, 490-496, 511, 513-519, 786 & J00-J06, J12-J18, J20-J22, J30, J31, J33, J34, J38, J39, J40-J47, J80-J86, J90-J92, J94, J96-J99, R04-R07, R09 \\
\hline
COPD & 491-492, 496 & J41-44 \\
\hline
Cardiovascular & 390-398, 401-417, 420-429, 440-449 & I00-I02, I05-I16, I20-I28, I30-I5A, I70-I79 \\
\hline
Psychiatric & 290-293, 295, 296, 298 & F06, F10-F16, F18-F21, F23, F25, F28-F33 \\
\hline
\end{tabular}
\caption{\textit{International Classification of Disease, Ninth Edition} (\textit{ICD}-9) and \textit{Tenth Edition} (\textit{ICD}-10) codes used in analysis, by medical condition. COPD, chronic obstructive pulmonary disease.} 
\label{tab:icd_codes}
\end{table}

\begin{table}[htbp]
    \includegraphics[width=\textwidth]{tables_figures/traditional_table1.png}
    \caption{Demographic characteristics of emergency department visits or hospitalizations included in the study in California from 2013--2019. Each observation represents a single visit, so individuals with multiple visits appear multiple times. COPD, chronic obstructive pulmonary disease.} 
    \label{traditional_table1}
\end{table}

\begin{table}[htbp]
    \includegraphics[width=.7\textwidth]{tables_figures/supp_exposure_table_same_day.png}
    \caption{Index days by same-day exposure status. Index days are days when a healthcare encounter for all-cause respiratory, COPD, cardiovascular, or psychiatric endpoints occurred. Wildfire \PM exposure represents the mean wildfire \PM concentration on that day. PSPS event magnitude (mild/moderate/severe) were classified using tertiles for both absolute and relative metrics. For our absolute (primary) metric, mild events were those affecting $\leq$ 61 people, moderate were those where the number of customers without power was $>$ 61 and $\leq$ 593, and severe was defined as events where $>$ 593 customers were without power. For our relative (secondary) metric, magnitude was classified as mild when our relative metric was $\leq$ 1, moderate when 1 $<$ relative metric $\geq$ 145, or severe when our relative metric was $>$ 145. COPD, chronic obstructive pulmonary disease; PSPS, public safety power shutoff}
    \label{supp_exposure_table_same_day}
\end{table}

\begin{table}[htbp]
    \includegraphics[width=.7\textwidth]{tables_figures/supp_exposure_table_lag0-3.png}
    \caption{Index days by \text{lag}$_{0\text{-}3}$ day exposure status. Index days are days when a healthcare encounter for all-cause respiratory, COPD, cardiovascular, or psychiatric endpoints occurred. Wildfire \PM exposure represents the mean wildfire \PM concentration for \text{lag}$_{0\text{-}3}$ days. PSPS event magnitude (mild/moderate/severe) were classified using tertiles for both absolute and relative metrics. For our absolute (primary) metric, mild events were those affecting $\leq$ 45 people, moderate were those where the number of customers without power was $>$ 45 and $\leq$ 458, and severe was defined as events where $>$ 458 customers were without power. For our relative (secondary) metric, magnitude was classified as mild when our relative metric was $\leq$ 1, moderate when 1 $<$ relative metric $\geq$ 145, or severe when our relative metric was $>$ 145. COPD, chronic obstructive pulmonary disease; PSPS, public safety power shutoff}
    \label{supp_exposure_table_lag0-3}
\end{table}


\begin{table}[htbp]
    \includegraphics[width=\textwidth]{tables_figures/results_table.png}
    \caption{Odds ratios and 95\% confidence intervals for the association between zip code-level severe PSPS events and all-cause respiratory-, COPD-, cardiovascular-, and psychiatric-related emergency department visits and hospital admissions in California from 2013 -- 2019. Estimates are from four separate conditional logistic regression models. All models include PSPS event exposure, \text{lag}$_{0\text{-}3}$ day wildfire \PM exposure, an interaction term between PSPS event and wildfire \PM exposure, and a spline on mean daily temperature. COPD, chronic obstructive pulmonary disease; PSPS, public safety power shutoff \newline \textit{*Note: The estimate for the interaction effect represents the additional multiplicative effect of joint exposure to wildfire \PM and PSPS event on the outcome of interest. It should not be interpreted as an odds ratio in isolation, but is rather a ratio of odds ratios.}}
    \label{supp_results_table}
\end{table}

\clearpage

\section{Supplemental figures.}

%%%% exposure figs 

\begin{figure}[htbp]
    \includegraphics[width=\textwidth]{tables_figures/seasonality_plot.pdf}
    \caption{Number of zip code-event days of PSPS events and mean wildfire \PM exposure, by month, in California during the study period (2013--2019). The number of zip code-event days of PSPS events reflects the sum over the 7-year period, while the mean wildfire \PM exposure represents means for each month over the 7-year study period. PSPS, public safety power shutoffs} 
    \label{seasonality_plot}
\end{figure}


\begin{figure}[h]
\includegraphics[width=.6\textwidth]{tables_figures/zips_included_map.pdf}    \caption{Spatial distribution of the \nzipcodesinanalysis zip codes in California included in this study. There were \nzipcodesresp, \nzipcodescopd, \nzipcodescardio, \nzipcodespsych \phantom{ } zip codes included for all-cause respiratory-, COPD-, cardiovascular-, and psychiatric-related emergency department visits and hospital admissions. COPD, chronic obstructive pulmonary disease.} 
    \label{zctas_included_map}
\end{figure}

%%%% results: mild

% absolute
\begin{figure}[htbp]
    \includegraphics[width=.9\textwidth]{tables_figures/age_stratified/supp_results_fig_mild_abs.pdf}
    \caption{Odds ratios and 95\% confidence intervals for the association between zip code-level mild PSPS events (absolute metric) and all-cause respiratory-, COPD-, cardiovascular-, and psychiatric-related emergency department visits and hospital admissions in California from 2013--2019 by age group. Estimates are from eight separate conditional logistic regression models, including one for each health endpoint -- exposure metric combination. All models include PSPS event exposure, \text{lag}$_{0\text{-}3}$ day wildfire \PM exposure, an interaction term between PSPS event and wildfire \PM exposure, and a spline on mean daily temperature. The absolute metric reflects the number of individuals exposed to a PSPS event, while the relative metric reflects the severity of the event. COPD, chronic obstructive pulmonary disease. \newline \textit{*Note: The estimate for the interaction effect represents the additional multiplicative effect of joint exposure to wildfire \PM and PSPS event on the outcome of interest. It should not be interpreted as an odds ratio in isolation, but is rather a ratio of odds ratios. PSPS, public safety power shutoff; COPD, chronic obstructive pulmonary disease}}
    \label{results_fig_20_64_abs}
\end{figure}

% relative
\begin{figure}[htbp]
    \includegraphics[width=.9\textwidth]{tables_figures/age_stratified/supp_results_fig_mild_hyb.pdf}
    \caption{Odds ratios and 95\% confidence intervals for the association between zip code-level mild PSPS events (relative metric) and all-cause respiratory-, COPD-, cardiovascular-, and psychiatric-related emergency department visits and hospital admissions in California from 2013--2019 by age group. Estimates are from eight separate conditional logistic regression models, including one for each health endpoint -- exposure metric combination. All models include PSPS event exposure, \text{lag}$_{0\text{-}3}$ day wildfire \PM exposure, an interaction term between PSPS event and wildfire \PM exposure, and a spline on mean daily temperature. The absolute metric reflects the number of individuals exposed to a PSPS event, while the relative metric reflects the severity of the event. COPD, chronic obstructive pulmonary disease. \newline \textit{*Note: The estimate for the interaction effect represents the additional multiplicative effect of joint exposure to wildfire \PM and PSPS event on the outcome of interest. It should not be interpreted as an odds ratio in isolation, but is rather a ratio of odds ratios. PSPS, public safety power shutoff; COPD, chronic obstructive pulmonary disease}}
    \label{results_fig_20_64_hyb}
\end{figure}


%%%% results: moderate
% absolute
\begin{figure}[htbp]
    \includegraphics[width=.9\textwidth]{tables_figures/age_stratified/supp_results_fig_mod_abs.pdf}
    \caption{Odds ratios and 95\% confidence intervals for the association between zip code-level moderate PSPS events (absolute metric) and all-cause respiratory-, COPD-, cardiovascular-, and psychiatric-related emergency department visits and hospital admissions in California from 2013--2019 by age group. Estimates are from eight separate conditional logistic regression models, including one for each health endpoint -- exposure metric combination. All models include PSPS event exposure, \text{lag}$_{0\text{-}3}$ day wildfire \PM exposure, an interaction term between PSPS event and wildfire \PM exposure, and a spline on mean daily temperature. The absolute metric reflects the number of individuals exposed to a PSPS event, while the relative metric reflects the severity of the event. COPD, chronic obstructive pulmonary disease. \newline \textit{*Note: The estimate for the interaction effect represents the additional multiplicative effect of joint exposure to wildfire \PM and PSPS event on the outcome of interest. It should not be interpreted as an odds ratio in isolation, but is rather a ratio of odds ratios. PSPS, public safety power shutoff; COPD, chronic obstructive pulmonary disease}}
    \label{results_fig_65_plus_abs}
\end{figure}

% relative 
\begin{figure}[htbp]
    \includegraphics[width=.9\textwidth]{tables_figures/age_stratified/supp_results_fig_mod_hyb.pdf}
    \caption{Odds ratios and 95\% confidence intervals for the association between zip code-level moderate PSPS events (relative metric) and all-cause respiratory-, COPD-, cardiovascular-, and psychiatric-related emergency department visits and hospital admissions in California from 2013--2019 by age group. Estimates are from eight separate conditional logistic regression models, including one for each health endpoint -- exposure metric combination. All models include PSPS event exposure, \text{lag}$_{0\text{-}3}$ day wildfire \PM exposure, an interaction term between PSPS event and wildfire \PM exposure, and a spline on mean daily temperature. The absolute metric reflects the number of individuals exposed to a PSPS event, while the relative metric reflects the severity of the event. COPD, chronic obstructive pulmonary disease. \newline \textit{*Note: The estimate for the interaction effect represents the additional multiplicative effect of joint exposure to wildfire \PM and PSPS event on the outcome of interest. It should not be interpreted as an odds ratio in isolation, but is rather a ratio of odds ratios. PSPS, public safety power shutoff; COPD, chronic obstructive pulmonary disease}}
    \label{results_fig_65_plus_hyb}
\end{figure}

%%%% results: severe
% absolute
\begin{figure}[htbp]
    \includegraphics[width=.9\textwidth]{tables_figures/age_stratified/supp_results_fig_sev_abs.pdf}
    \caption{Odds ratios and 95\% confidence intervals for the association between zip code-level severe PSPS events (absolute metric) and all-cause respiratory-, COPD-, cardiovascular-, and psychiatric-related emergency department visits and hospital admissions in California from 2013--2019 by age group. Estimates are from eight separate conditional logistic regression models, including one for each health endpoint -- exposure metric combination. All models include PSPS event exposure, \text{lag}$_{0\text{-}3}$ day wildfire \PM exposure, an interaction term between PSPS event and wildfire \PM exposure, and a spline on mean daily temperature. The absolute metric reflects the number of individuals exposed to a PSPS event, while the relative metric reflects the severity of the event. COPD, chronic obstructive pulmonary disease. \newline \textit{*Note: The estimate for the interaction effect represents the additional multiplicative effect of joint exposure to wildfire \PM and PSPS event on the outcome of interest. It should not be interpreted as an odds ratio in isolation, but is rather a ratio of odds ratios. PSPS, public safety power shutoff; COPD, chronic obstructive pulmonary disease}}
    \label{results_fig_65_plus_abs}
\end{figure}

% relative
\begin{figure}[htbp]
    \includegraphics[width=.9\textwidth]{tables_figures/age_stratified/supp_results_fig_sev_hyb.pdf}
    \caption{Odds ratios and 95\% confidence intervals for the association between zip code-level severe PSPS events (relative metric) and all-cause respiratory-, COPD-, cardiovascular-, and psychiatric-related emergency department visits and hospital admissions in California from 2013--2019 by age group. Estimates are from eight separate conditional logistic regression models, including one for each health endpoint -- exposure metric combination. All models include PSPS event exposure, \text{lag}$_{0\text{-}3}$ day wildfire \PM exposure, an interaction term between PSPS event and wildfire \PM exposure, and a spline on mean daily temperature. The absolute metric reflects the number of individuals exposed to a PSPS event, while the relative metric reflects the severity of the event. COPD, chronic obstructive pulmonary disease. \newline \textit{*Note: The estimate for the interaction effect represents the additional multiplicative effect of joint exposure to wildfire \PM and PSPS event on the outcome of interest. It should not be interpreted as an odds ratio in isolation, but is rather a ratio of odds ratios. PSPS, public safety power shutoff; COPD, chronic obstructive pulmonary disease}}
    \label{results_fig_65_plus_hyb}
\end{figure}



%%%% sensitivity analysis: 1 week duration

\begin{figure}[htbp]
    \includegraphics[width=.9\textwidth]{tables_figures/1week_duration/results_fig_1week.pdf}
    \caption{Odds ratios and 95\% confidence intervals for the association between zip code-level PSPS event duration and all-cause respiratory-, COPD-, cardiovascular-, and psychiatric-related emergency department visits and hospital admissions in California from 2013--2019 among those 20 years of age and older (full study population). Estimates are from four separate conditional logistic regression models -- one for each health endpoint. All models include PSPS duration, lag 0-6 day wildfire \PM exposure, an interaction term between PSPS duration and wildfire \PM exposure, and a spline on mean daily temperature. \newline \textit{*Note: The estimate for the interaction effect represents the additional multiplicative effect of joint exposure to wildfire \PM and PSPS event on the outcome of interest. It should not be interpreted as an odds ratio in isolation, but is rather a ratio of odds ratios. PSPS, public safety power shutoff; COPD, chronic obstructive pulmonary disease}}
    \label{1week_duration_fig}
\end{figure}

\begin{figure}[htbp]
    \includegraphics[width=.9\textwidth]{tables_figures/1week_duration/results_fig_1week_20_64.pdf}
    \caption{Odds ratios and 95\% confidence intervals for the association between zip code-level PSPS event duration and all-cause respiratory-, COPD-, cardiovascular-, and psychiatric-related emergency department visits and hospital admissions in California from 2013--2019 among those ages 20-64. Estimates are from four separate conditional logistic regression models -- one for each health endpoint. All models include PSPS duration, lag 0-6 day wildfire \PM exposure, an interaction term between PSPS duration and wildfire \PM exposure, and a spline on mean daily temperature. \newline \textit{*Note: The estimate for the interaction effect represents the additional multiplicative effect of joint exposure to wildfire \PM and PSPS event on the outcome of interest. It should not be interpreted as an odds ratio in isolation, but is rather a ratio of odds ratios. PSPS, public safety power shutoff; COPD, chronic obstructive pulmonary disease}}
    \label{1week_duration_fig_young}
\end{figure}

\begin{figure}[htbp]
    \includegraphics[width=.9\textwidth]{tables_figures/1week_duration/results_fig_1week_65_plus.pdf}
    \caption{Odds ratios and 95\% confidence intervals for the association between zip code-level PSPS event duration and all-cause respiratory-, COPD-, cardiovascular-, and psychiatric-related emergency department visits and hospital admissions in California from 2013--2019 among those 65 years of age and older. Estimates are from four separate conditional logistic regression models -- one for each health endpoint. All models include PSPS duration, lag 0-6 day wildfire \PM exposure, an interaction term between PSPS duration and wildfire \PM exposure, and a spline on mean daily temperature. \newline \textit{*Note: The estimate for the interaction effect represents the additional multiplicative effect of joint exposure to wildfire \PM and PSPS event on the outcome of interest. It should not be interpreted as an odds ratio in isolation, but is rather a ratio of odds ratios. PSPS, public safety power shutoff; COPD, chronic obstructive pulmonary disease}}
    \label{1week_duration_fig_old}
\end{figure}



\clearpage
\end{document}